% 原书第 96—100 页；6.4。
\section{反事实}
\label{sec:ch06-counterfactuals}

\textbf{反事实（counterfactual）}的定义与解释，在文献中受到广泛关注。它所处理的情形如下：假设你正在打扑克，起手牌是 $\clubsuit\mathrm J$ 和 $\clubsuit3$（这有时被戏称为“伐木工”，lumberjack，取 tree 与 three 的谐音，再加上 jack）。你估计获胜概率太小，不想再输钱，于是弃牌（fold）。接下来翻开三张公共牌，即翻牌（flop）：$\clubsuit4$、$\clubsuit\mathrm Q$ 和 $\clubsuit2$。你会作出一种典型的反事实反应：“如果当时继续玩，我的胜算就很大了。”（五张同花色牌称为“同花”，在牌型中位列第五；甚至还有机会组成位列第二的“同花顺”。）这一陈述将观测数据——手牌与翻牌——纳入模型，再分析一种干预分布：让你继续游戏，而环境的其余部分保持不变，即牌面相同。形式上，这相当于先通过条件化更新结构因果模型中的噪声分布，再实施干预。

\begin{definition}[反事实]
考虑节点为 $\mathbf X$ 的结构因果模型 $\mathfrak C:=(\mathbf S,P_{\mathbf N})$。给定某些观测值 $\mathbf x$，通过替换噪声变量的分布，定义反事实结构因果模型：
\[
 \mathfrak C_{\mathbf X=\mathbf x}
 :=\bigl(\mathbf S,P_{\mathbf N}^{\mathfrak C\mid\mathbf X=\mathbf x}\bigr),
\]
其中 $P_{\mathbf N}^{\mathfrak C\mid\mathbf X=\mathbf x}:=P_{\mathbf N\mid\mathbf X=\mathbf x}$。\footnote{在连续情形下，这一定义会遇到测度论问题，因为条件分布通常只在相差零测集的意义下确定。为简化处理，我们将反事实限制在离散情形，即噪声分布具有概率质量函数的情形。若变量连续且具有密度，则不以 $\mathbf X=\mathbf x$ 为条件，而以满足 $P(\mathbf X\in A)>0$ 的事件 $\mathbf X\in A$ 为条件。}
新模型中的噪声变量不再需要相互独立。反事实陈述，现在可以视为新的反事实结构因果模型中的 do 陈述。
\end{definition}

这一定义可以推广到只观测部分变量，而非完整向量 $\mathbf X=\mathbf x$ 的情形。

\begin{example}[计算反事实]
考虑以下结构因果模型：
\[
 \begin{aligned}
 X&:=N_X\\
 Y&:=X^2+N_Y\\
 Z&:=2\cdot Y+X+N_Z,
 \end{aligned}
\]
其中 $N_X,N_Y,N_Z\overset{\mathrm{iid}}{\sim}\mathrm U(\{-5,-4,\ldots,4,5\})$，均匀分布在 $-5$ 到 $5$ 的整数上。现在，假设观测到 $(X,Y,Z)=(1,2,4)$。由于此时可以从观测值唯一重建全部噪声项，$P_{\mathbf N}^{\mathfrak C\mid\mathbf X=\mathbf x}$ 将全部概率质量放在 $(N_X,N_Y,N_Z)=(1,1,-1)$ 上。因此，在观测 $(X,Y,Z)=(1,2,4)$ 的背景下，有反事实陈述：“如果 $X$ 当时为 $2$，$Z$ 就会为 $11$。”本书将这种句子解释为：“如果当时将 $X$ 设为 $2$，$Z$ 就会为 $11$。”用数学语言来说，$P_Z^{\mathfrak C\mid\mathbf X=\mathbf x;\dox{X:=2}}$ 将全部概率质量放在 $11$ 上。同样，可以得到：“如果 $X$ 当时为 $2$，$Y$ 就会为 $5$”，以及“如果 $Y$ 当时为 $5$，$Z$ 就会为 $10$”。
\end{example}

构造反事实涉及多个步骤，因此记号看起来相当复杂。\footnote{Pearl［2009］采用较简洁的记法 $Z_y(u)$，其中下标 $y$ 表示干预 $\dox{Y:=y}$；$u$ 表示有关误差项的额外信息，例如由 $\mathbf X=\mathbf x$ 所蕴含的信息。他将误差项称为 $u$。}
希望下图有助于进一步说明。
% 原书第 97 页的无编号记号说明图。
\begin{center}
\begin{tikzpicture}[causal,x=1.1cm,y=1cm]
 \node[font=\large] (ch06cf-formula) at (0,0)
 {$P_Z^{\mathfrak C\mid\mathbf X=\mathbf x;\dox{Y:=2}}$};
 \node[align=center,font=\small] (ch06cf-one) at (-3.15,1.1) {1. 出发模型 $\mathfrak C$};
 \node[align=center,font=\small] (ch06cf-two) at (-.4,1.9) {2. 观测数据 $\mathbf X=\mathbf x$};
 \node[align=center,font=\small] (ch06cf-three) at (3.1,1.1) {3. 干预 $\dox{Y:=2}$};
 \node[align=center,font=\small] (ch06cf-four) at (1.8,-1) {4. 关注的变量 $Z$};
 \draw[ed] (ch06cf-one.south)--(-.72,.2);
 \draw[ed] (ch06cf-two.south)--(-.08,.3);
 \draw[ed] (ch06cf-three.south)--(.9,.3);
 \draw[ed] (ch06cf-four.north)--(-.78,-.2);
\end{tikzpicture}
\end{center}


反事实陈述高度依赖于结构因果模型的结构。例 6.19 给出两个模型：它们诱导出相同的图、观测分布和干预分布，却导出不同的反事实陈述。后文会称它们为“概率等价且干预等价”，但“反事实不等价”（见定义 6.47）。

\begin{example}
设 $N_1,N_2\sim\operatorname{Ber}(0.5)$，$N_3\sim\mathrm U(\{0,1,2\})$，且三个变量相互独立。也就是说，$N_1,N_2$ 服从参数为 $0.5$ 的 Bernoulli 分布，$N_3$ 在 $\{0,1,2\}$ 上均匀分布。定义两个不同的结构因果模型。首先考虑 $\mathfrak C_A$：
\[
 \begin{aligned}
 X_1&:=N_1\\
 X_2&:=N_2\\
 X_3&:=(\mathbf1_{N_3>0}\cdot X_1+\mathbf1_{N_3=0}\cdot X_2)
 \cdot\mathbf1_{X_1\ne X_2}+N_3\cdot\mathbf1_{X_1=X_2}.
 \end{aligned}
\]
如果 $X_1$ 与 $X_2$ 的值不同，就根据 $N_3$，选择 $X_3=X_1$ 或 $X_3=X_2$；否则令 $X_3=N_3$。第二个模型 $\mathfrak C_B$ 只在后一种情形下与 $\mathfrak C_A$ 不同：
\[
 \begin{aligned}
 X_1&:=N_1\\
 X_2&:=N_2\\
 X_3&:=(\mathbf1_{N_3>0}\cdot X_1+\mathbf1_{N_3=0}\cdot X_2)
 \cdot\mathbf1_{X_1\ne X_2}+(2-N_3)\cdot\mathbf1_{X_1=X_2}.
 \end{aligned}
\]
两个模型导出相同的观测分布；对于任意可能的干预，它们也导出相同的干预分布。\footnote{本例的观测分布相对于其背后的图 $X_1\to X_3\leftarrow X_2$ 满足因果极小性；见定义 6.33。第 3.4 节给出了另一个更简单的例子，但该例违反因果极小性。}
不过，二者在一项反事实陈述上不同。假设我们已经观测到 $(X_1,X_2,X_3)=(1,0,0)$，并关心反事实问题：“如果 $X_1$ 当时为 $0$，$X_3$ 会是多少？”两个模型都推出 $N_3=0$。因此，当反事实地改变 $X_1$ 时，$\mathfrak C_A$ 与 $\mathfrak C_B$ 对 $X_3$“预测”的取值不同，分别为 $0$ 和 $2$。
\end{example}

上例有两方面含义。（1）这两个结构因果模型对应同一个因果图模型（见第 6.5.2 节）。在这个意义上，因果图模型的表达能力不足以预测反事实。（2）第 6.8 节将把干预分布与现实世界中的随机化实验联系起来。在本例中，随机化试验和观测数据都无法区分 $\mathfrak C_A$ 与 $\mathfrak C_B$。因此，如果关心反事实陈述，就需要能够区分二者的额外假设。

下面总结反事实的一些性质。

\begin{remark}
\begin{enumerate}[label=(\roman*)]
 \item 反事实陈述不具有传递性。在例 6.18 中，给定观测 $(X,Y,Z)=(1,2,4)$，我们得到：
 \begin{quote}
 “如果 $X$ 当时为 $2$，$Y$ 就会为 $5$”；\\
 “如果 $Y$ 当时为 $5$，$Z$ 就会为 $10$”；\\
 “如果 $X$ 当时为 $2$，$Z$ 就\emph{不会}为 $10$”。
 \end{quote}
 因此，不能简单地引入新变量，例如 $\widetilde X$ 和 $\widetilde Y$，再将“如果 $X$ 当时为 $2$，$Y$ 就会为 $5$”解释为“$\widetilde X=2\Rightarrow\widetilde Y=5$”这样的逻辑蕴含。在上例中，不具有传递性，是因为从 $X$ 到 $Z$ 存在一条直接联系，即存在一条不经过 $Y$ 的路径。干预分布也有类似的反例。

 \item 人们经常以反事实方式思考。“我当时应该坐火车。”“还记得我们 2000 年 9 月 11 日飞往纽约的那趟航班吗？想象一下，如果晚一年坐那趟航班，会怎样！”“我们应该在 2014 年 12 月投资瑞士法郎！”这些只是几个例子。有意思的是，有时即使依据当时掌握的信息，我们已经作出了最优决策，仍会这样想。假设有人提出，只要你猜中抛硬币的结果，就给你 10,000 美元。你猜正面，却输了。有些人接着会想：“为什么我当时不说反面？”尽管根本不可能事先知道结果。Roese［1997］、Byrne［2007］等讨论了反事实思维的心理学意义。反事实陈述是否包含有助于今后作出更好决策的信息，是一个有趣的问题，但超出本书范围；另见 Pearl［2009，第 4 章］。

 \item 本书也不讨论反事实在法律制度中的作用。是否应当以及如何将反事实作为裁判依据，是一个值得研究的问题（见例 3.4）。

 \item 人类思考反事实由来已久；它是历史学家常用的工具。例如，Titus Livius 在公元前 25 年讨论过：如果 Alexander the Great 没有死在亚洲，而是进攻罗马，会发生什么［Geradin and Girgenson, 2011］。Paul 的《哥林多前书》（7:29—7:31）写道：“弟兄们，我是说，时候短促了。从今以后，有妻子的，要像没有妻子；哀哭的，要像不哀哭；快乐的，要像不快乐；置买的，要像无有所得；用世物的，要像不用世物到尽头。”

 \item 可以将干预陈述理解为（随机化）实验的数学构造。对于反事实陈述，现实世界中没有类似的对应关系。人们可能会推测，许多反事实陈述无法被否证，因此不应在科学研究中使用［参见 Popper, 2002］。不过，有时我们确实可以作出可否证的反事实陈述，例如样本中某个个体的噪声项实际取值在事后显现出来时；见例 3.4。此外，上面讨论的反事实，是设定某个结构因果模型之后的推论。因此，否证的对象也可以是结构因果模型，而非某个特定的反事实陈述。例如，借助该模型所指向的科学领域中的方法，有时可以做到，有时则不行。\footnote{注意，第 3.4 节所述的重参数化自由度始终存在。}
\end{enumerate}
\end{remark}

上述说明可以作为进一步思考的起点。我们不再深入讨论如何解释反事实陈述，以及应当或能够如何将其用于法庭案件等问题。这些思考中有许多超出了我们的专业范围。读者可参阅 Halpern［2016］，其中讨论了某个事件是另一个事件的“实际原因”（actual cause）究竟意味着什么。
\FloatBarrier
